\documentclass[11pt,a4paper]{article}

% ------------------------------------------------------------------
%  Documentación técnica \textemdash{} Plataforma de Monitoreo Knop Laboratorios
% ------------------------------------------------------------------
\usepackage[utf8]{inputenc}
\usepackage[T1]{fontenc}
\usepackage{lmodern}
\usepackage{amsmath}
\usepackage[spanish,es-noindentfirst]{babel}
\usepackage[a4paper,margin=2.4cm]{geometry}
\usepackage{booktabs}
\usepackage{longtable}
\usepackage{array}
\usepackage{xcolor}
\usepackage{fancyhdr}
\usepackage{enumitem}
\usepackage{titlesec}
\usepackage[hidelinks]{hyperref}

\definecolor{knopgreen}{HTML}{09612D}
\definecolor{knopink}{HTML}{0F1C14}
\definecolor{codebg}{HTML}{F4F7F4}

\titleformat{\section}{\Large\bfseries\color{knopgreen}}{\thesection}{0.6em}{}
\titleformat{\subsection}{\large\bfseries\color{knopink}}{\thesubsection}{0.6em}{}

\newcommand{\code}[1]{\texttt{\small #1}}

\pagestyle{fancy}
\fancyhf{}
\lhead{\textcolor{knopgreen}{\textbf{Knop Laboratorios}} · Monitoreo ambiental}
\rhead{\thepage}
\renewcommand{\headrulewidth}{0.4pt}

\title{\vspace{-1.2cm}\textbf{\textcolor{knopgreen}{Plataforma de Monitoreo Ambiental}}\\[2pt]
\large Documentación técnica del sistema}
\author{Knop Laboratorios \textemdash{} Ingeniería y Mantención}
\date{Junio 2026 · v1.1.1}

\begin{document}
\shorthandoff{"<>}
\maketitle
\thispagestyle{fancy}

\begin{abstract}
\noindent Este documento describe la plataforma web de monitoreo ambiental construida para
Knop Laboratorios: su arquitectura, qué datos entrega la API del proveedor (sensores LoRaWAN
de diferencial de presión y termohigrómetros), qué información extraemos y calculamos, y qué
funcionalidades ofrece hoy y puede ofrecer a futuro.
\end{abstract}

\tableofcontents
\vspace{0.6cm}

% ==================================================================
\section{Cómo funciona el sistema}
% ==================================================================

La plataforma es una aplicación web (Next.js + TypeScript, desplegada en Netlify) que
\textbf{no almacena datos propios}: consume en vivo las API del sistema de monitoreo del
proveedor (Softronica) sobre los sensores instalados en las áreas controladas de Knop.

\subsection{Flujo de datos}

\begin{enumerate}[leftmargin=1.4em]
  \item El \textbf{navegador} (interfaz React) solicita datos a nuestras rutas internas
        \code{/api/knop/*}.
  \item Esas rutas actúan como \textbf{proxy}: en el servidor consultan la API del proveedor,
        \textbf{normalizan} los tipos (números, fechas, codificación UTF-8) y devuelven JSON limpio.
  \item La API del proveedor (\code{PHP}) lee la base de datos alimentada por los
        \textbf{sensores LoRaWAN} y responde en JSON.
\end{enumerate}

\begin{center}
\fbox{\parbox{0.92\textwidth}{\centering
Navegador (React) \;$\rightarrow$\; \code{/api/knop/*} (proxy + normalización)
\;$\rightarrow$\; API PHP del proveedor \;$\rightarrow$\; Sensores LoRaWAN
}}
\end{center}

\subsection{¿Por qué un proxy?}
\begin{itemize}[leftmargin=1.4em]
  \item \textbf{Tipado y limpieza}: la API original entrega números como texto y con codificación
        variable; el proxy los normaliza.
  \item \textbf{Desacople}: si en el futuro Knop cambia de proveedor o monta su propio backend,
        solo cambia el proxy, no la interfaz.
  \item \textbf{Control de caché}: evita que datos de un sensor se mezclen con otro
        (ver Sección~\ref{sec:notas}).
\end{itemize}

\subsection{Vistas de la aplicación}
\begin{center}
\begin{tabular}{@{}p{4.6cm}p{9.8cm}@{}}
\toprule
\textbf{Ruta} & \textbf{Descripción} \\
\midrule
\code{/} & Inicio / acceso a las vistas. \\
\code{/monitoreo/sdp} & Diferencial de presión (Pa / inH$_2$O) con rango operacional. \\
\code{/monitoreo/termohigrometros} & Temperatura y humedad con rangos de aceptación. \\
\code{/informe} & Informe estadístico (presión \emph{o} termohigrómetro). \\
\bottomrule
\end{tabular}
\end{center}

% ==================================================================
\section{Qué entrega la API del proveedor}
% ==================================================================

Base: \code{https://newenergy.softronica.cl/knop/monitoreo/}. Todos los endpoints responden
JSON y son de lectura. Hay dos familias de sensores: \textbf{Diferencial de Presión (DP)} y
\textbf{Termohigrómetros (STH)}.

\subsection{Endpoints}
\begin{center}
\begin{longtable}{@{}p{4.1cm}p{4.0cm}p{6.3cm}@{}}
\toprule
\textbf{Endpoint} & \textbf{Parámetros} & \textbf{Devuelve} \\
\midrule
\endhead
\code{sdp/deviceDP.php} & \textemdash{} & Lista de sensores de presión \\
\code{sdp/kpiDP.php} & \code{devEui}, \code{preset} ó \code{start+end}, \code{agg} & Serie temporal de presión \\
\code{sdp/rangoDP.php} & \code{devEui} & Rango operacional de presión \\
\code{device.php} & \textemdash{} & Lista de termohigrómetros \\
\code{kpiSTH.php} & \code{deviceName}, \code{preset} ó \code{start+end}, \code{agg} & Serie temporal de temp./humedad \\
\code{rangoSTH.php} & \code{deviceName} & Rango de aceptación de temp./humedad \\
\bottomrule
\end{longtable}
\end{center}

\subsection{Campos de los datos}

\paragraph{Sensor de presión (\code{deviceDP.php}).}
\code{devEui}, \code{identificador}, \code{ubicacion}, \code{area}, \code{seccion},
\code{tipo\_rango}, \code{label}.

\paragraph{Serie de presión (\code{kpiDP.php}).}
\begin{center}
\begin{tabular}{@{}p{4.8cm}p{9.6cm}@{}}
\toprule
\textbf{Campo} & \textbf{Significado} \\
\midrule
\code{bucket\_time} & Marca de tiempo del intervalo agregado (eje X). \\
\code{last\_datetime} & Fecha/hora de la última lectura del intervalo. \\
\code{Differential\_pressure\_Pa} & Presión diferencial en Pascales. \\
\bottomrule
\end{tabular}
\end{center}
\noindent\textit{Nota: la serie de presión \textbf{no} incluye batería.}

\paragraph{Serie de termohigrómetro (\code{kpiSTH.php}).}
\begin{center}
\begin{tabular}{@{}p{3.4cm}p{11.0cm}@{}}
\toprule
\textbf{Campo} & \textbf{Significado} \\
\midrule
\code{bucket\_time} & Marca de tiempo del intervalo agregado. \\
\code{tempC\_SHT} & Temperatura ($^\circ$C) del sensor principal (SHT). \\
\code{hum\_SHT} & Humedad relativa (\%). \\
\code{batV} & \textbf{Voltaje de batería del sensor} ($\approx$3.6 V). \\
\code{tempC\_DS} & Temperatura de una segunda sonda (DS); no se muestra aún. \\
\bottomrule
\end{tabular}
\end{center}
\noindent\textbf{Sobre la batería:} el dato \code{batV} \emph{proviene directamente de la API}
de termohigrómetros, no es estimado. Solo está disponible para STH; los sensores de presión
no exponen batería.

\paragraph{Rangos (\code{rangoDP.php} / \code{rangoSTH.php}).}
DP: \code{min\_pa}, \code{max\_pa}, \code{tipo\_rango\_dp}, \code{descripcion\_rango}.\\
STH: \code{temp\_min}, \code{temp\_max}, \code{hum\_min}, \code{hum\_max}, \code{tipo\_rango\_sth}.

\subsection{Ejemplo de respuesta (termohigrómetro)}
\begin{verbatim}
[{"bucket_time":"2026-06-22 21:50:00","tempC_SHT":"20.9",
  "hum_SHT":"38.4","batV":"3.63","tempC_DS":"2.09"}, ...]
\end{verbatim}

% ==================================================================
\section{Qué extraemos y calculamos}
% ==================================================================

\subsection{Reglas del sistema (reutilizadas del original)}
\begin{itemize}[leftmargin=1.4em]
  \item \textbf{Conversión de presión}: $1\ \text{inH}_2\text{O} = 249.08891\ \text{Pa}$.
  \item \textbf{Agregación por periodo} (minutos por punto):
        24h$\to$5, 2d$\to$5, 3d$\to$5, 7d$\to$15, 30d$\to$60, 12m$\to$1440.
  \item En rango personalizado, la agregación se estima según la cantidad de días.
\end{itemize}

\subsection{Métricas derivadas (informe estadístico)}
A partir de la serie de cada sensor y su rango, calculamos:
\begin{center}
\begin{tabular}{@{}p{4.8cm}p{9.6cm}@{}}
\toprule
\textbf{Métrica} & \textbf{Cómo se obtiene} \\
\midrule
Mínimo / Promedio / Máximo & Sobre las lecturas del periodo. \\
\% de cumplimiento & Proporción de lecturas dentro del rango de aceptación. \\
Tiempo fuera de rango & N.\textsuperscript{o} de intervalos fuera $\times$ minutos por intervalo. \\
Tendencia & Comparación del último tercio vs. el tercio previo. \\
Estado general & Normal / Advertencia / Alerta según la última lectura vs.\ rango. \\
Alarmas con histéresis & Se activan al cruzar el umbral y se reponen al volver dentro del 5\% del rango. \\
\bottomrule
\end{tabular}
\end{center}

\noindent En termohigrómetros estas métricas se calculan \textbf{por separado para temperatura
y humedad}, y el estado general es el peor de ambas.

% ==================================================================
\section{Qué se puede hacer}
% ==================================================================

\subsection{Funcionalidades actuales}
\begin{itemize}[leftmargin=1.4em]
  \item Réplica mejorada de las páginas de monitoreo (presión y termohigrómetros) con
        identidad visual Knop y bandas de rango siempre visibles.
  \item Selección de sensor y de periodo (24h a 12 meses) y rango de fechas personalizado.
  \item Auto-actualización en vivo (cada 60 s) y exportación a \textbf{Excel}.
  \item \textbf{Informe estadístico} (presión o termohigrómetro) con estado, cumplimiento,
        tiempo fuera de rango, estadísticas, alarmas y exportación a \textbf{PDF}.
\end{itemize}

\subsection{Extensiones posibles}
\begin{itemize}[leftmargin=1.4em]
  \item \textbf{Vista por área/sección} y \textbf{comparación} entre varios sensores a la vez.
  \item \textbf{Panel general} (todos los sensores en una grilla con su estado).
  \item Aprovechar \code{tempC\_DS} (segunda sonda) y \textbf{alertas de batería baja} (usando \code{batV}).
  \item \textbf{Notificaciones} (correo / mensajería) ante alarmas o sensores sin reportar.
  \item \textbf{Informes PDF programados} (p.\,ej.\ mensuales) para soporte de calidad/GMP.
  \item \textbf{Backend propio} con histórico, para no depender del proveedor.
  \item \textbf{Autenticación} y dominio personalizado.
\end{itemize}

% ==================================================================
\section{Notas operativas}\label{sec:notas}
% ==================================================================
\begin{itemize}[leftmargin=1.4em]
  \item \code{kpiSTH.php} requiere el \emph{nombre completo} del sensor
        (\code{"CÓDIGO - Ubicación"}), no solo el código.
  \item Las rutas del proxy responden \code{Cache-Control: no-store} para evitar que el CDN
        mezcle datos entre sensores; el cacheo del origen se hace en el servidor.
  \item Datos consumidos con autorización del cliente. El proxy oculta la URL del proveedor.
\end{itemize}

\vfill
\noindent\textcolor{knopgreen}{\rule{\textwidth}{0.8pt}}\\[2pt]
{\footnotesize Repositorio: \url{https://github.com/FranJGT/monitoreo-knop} \quad·\quad
Producción: \url{https://superlative-truffle-37d616.netlify.app}}

\end{document}
